\subsection{Using the translate mean in the window construction}\label{red:translate}

Fix the character, \(\varepsilon\), the common host at \(x\), and the
weights and dimensions in part 2 of Formal input~\ref{red:design}.
Write \(\mathcal A_u=\sum_{t,j}\mathrm{ASK}_{u,t,j}\),
\(\mathcal M_u=\sum_{t,j}\mathrm{MAIN}_{u,t,j}\),
\(D_u=\mathcal A_u-\mathcal M_u\). The expectation throughout this
subsection is uniform on the subgroup \(A_0\) in
Proposition~\ref{ar:mean}. The original design \code{topDesWin} selects
one family. Its later application in \code{ch\_top} bounds each tied translate. To apply a finite mean lemma, first restrict the tied data on the host group to \(A_0\).
The source file also gives a
\src{src/Formalism/SizeWindowBrackets/ChainSlopeWindows/LevelOfCongruence.lean}{199}{mean form of the slope comparison}.
This alternative takes two mean loss bounds. Its proof first obtains
a geometric inequality for every translate separately. Its text is
present in the pinned source, while the saved classical dependency
walk uses the pointwise declaration on line 175. We give the scalar
deduction and then replace those separate HN comparisons by
Corollary~\ref{trmean:budget}. When the host group is larger than
\(A_0\), a bound for its whole-group mean alone gives no bound for
the mean over this particular subgroup.

Let \(G_u\) be the recorded divisibility gain and \(I_u\) the
local orbit loss (the quantities \(\mathrm{Ga}\) and \(\mathrm{IO}\)
in the slope comparison). Let \(T_{ui}>0\) be the word level.
Use \(a_E,\mathcal L_a,\mathcal L,b',X''\) for the numerical
parameters \(a,L_a,L,b',X''\) of the common comparison and
\(c_s=c_{\mathrm{Sl1}}\ge0\). Put \(A=\sum_t a_t\) and
\[
\begin{aligned}
 C_N&=c_{\rm count}(\mathcal L_a+m^2\mathcal L_a/k),&
 C_M&=c_{\rm slope} n_Ta_E\mathcal L_a(1+m/k),\\
 E_{\rm div}&=c_{\rm div}\big(n_Ta_E\mathcal L_a+
                       m(B_0+N_c)x^{e_{\rm div}\varepsilon}\big),&
 K_0&=A C_N+E_{\rm div}+mc_EB_0+C_M.
\end{aligned}
\]
The word loss and its rank mean are
\[
 \ell_{ui}=c_s I_u+c_s n_Tm(a_E/k+b')\log T_{ui}+c_sX''+1,
 \qquad \ell_u=\sum_i r_{ui}\ell_{ui}/r_u .
\]
The counting bound \(\log T_{ui}\le c_{\rm count}m\mathcal L\) gives
\(\ell_u\le c_s I_u+K_7\), where
\(K_7=c_s n_Tm(a_E/k+b')c_{\rm count}m\mathcal L+c_sX''+1\).
Here and below the size parameters multiplying this counting bound
are nonnegative.

\begin{proposition}[A sufficient mean comparison at one scale]\label{trmean:assembly}
Retain the identification of the descended columns, their original
shifted-source slopes and the target slopes of Section~\ref{geo:translate}
with the data just chosen. Suppose the division and lower slope
comparisons on these data give
\begin{equation}\label{trmean:lower}
 \mathbb E(G_u-\mathcal M_u)
       \le \mathbb E(s_u-\mu_u)+K_0+\mathbb E\ell_u .
\end{equation}
If \(F_x\) is either right side of \eqref{trmean:gap} or
\eqref{trmean:budget-eq} under its stated hypotheses, then
\begin{equation}\label{trmean:test}
 \mathbb E D_u\le K_0+K_7+F_x+
                  \mathbb E(\mathcal A_u-G_u+c_s I_u).
\end{equation}
In particular, bounding this right side by \(E_x(\mathcal D_x)\) in
Section~\ref{scope:account} proves \eqref{scope:average} at that scale.
\end{proposition}
\begin{proof}
To see precisely what supplies \eqref{trmean:lower} in the original
word estimates, write \(W_u=\sum_j f^{-1}\sum_\rho\log w_{u,j,\rho}\)
and \(\mathcal M_u=S_0-AW_u\) by the definition of the centre.
Here \(S_0\) is independent of \(u\). Write \(v_0>0\) for the
archimedean scale in the count and \(\ell_g\) for its additive height
parameter. Set
\[
 S_D=S_0-mA\log(\pi/v_0),\qquad
 \Lambda_u=W_u+m\log(v_0/\pi)+m\ell_g-C_N .
\]
Expansion gives \(S_D-A\Lambda_u+mA\ell_g=\mathcal M_u+AC_N\).
For a word let \(F_u^{\,0}\) be the recorded shift common to its
translate, \(E_{ui}\) the sum of its local exponent corrections,
and \(\mu^M_{ui}\) its base lattice slope. The input inequalities are
\[
\begin{aligned}
 \mu_u+h_{ui}-F_u^{\,0}
    &\le S_D-A\log T_{ui}+E_{\rm div},&
 \log T_{ui}&\ge\Lambda_u,\\
 \mu(S_{ui})&\ge\mu^M_{ui}+G_u+F_u^{\,0}-E_{ui}-\ell_{ui},&
 \mu^M_{ui}&\ge-C_M,\\
 E_{ui}&\le m(A\ell_g+c_EB_0).&&
\end{aligned}
\]
Insert the lower bound on \(F_u^{\,0}-h_{ui}\) from the first line
into the second. Since \(A\ge0\), one may replace \(A\log T_{ui}\)
by \(A\Lambda_u\) in the resulting lower bound. The displayed identity
then gives \(G_u-\mathcal M_u\le\mu(S_{ui})-h_{ui}-\mu_u+
K_0+\ell_{ui}\). Averaging by word rank and then by translate proves
\eqref{trmean:lower}. An independent proof of just its mean would
also suffice. Finally use
\(D_u=(\mathcal A_u-G_u)+(G_u-\mathcal M_u)\), the word counting
bound, and the single slope comparison. This gives \eqref{trmean:test}.
\end{proof}

\begin{corollary}[The original bracket from one signed loss]\label{trmean:loss}
Use the common numerical parameters \((\delta_*,\eta_*,e_*)\) of
the original error bracket, including its large-\(x\) size and
error-window hypotheses. Assume the mass and compression bounds of
Corollary~\ref{trmean:budget} with this \(\delta_*\) and the mean
frame bounds \(e\le e_*,\ \eta\le\eta_*\). Write \(X=x^{c_{\rm HH}}\),
\(H=Xb\), \(b>0\). Suppose \(\mathcal J=T_{\rm out}=0\) and
\(U\le K_3+\mathbb E\ell_u\), with \(K_3\) the variation bound below.
In addition to \eqref{trmean:lower} assume
\begin{equation}\label{trmean:signed}
 \mathbb E(\mathcal A_u-G_u+\kappa I_u)\le(1+\kappa)L_P,\qquad
 \kappa=(1+C_\phi)c_s,\quad
 C_\phi=\frac{4\log(1+X)}{1-\delta_*},
\end{equation}
where \(L_P=c_{\rm loss}n_Tmb\,x^{e_{\rm loss}\varepsilon}\).
Then the original numerical mean bracket bounds \(\mathbb E D_u\).
In particular the same numerical size hypotheses as in Formal
input~\ref{red:design} give the right side of \eqref{scope:average}.
\end{corollary}
\begin{proof}
For \(a\ge0\), \(a\log(1+Xb/a)\le(a+b)\log(1+X)\), with value
zero at \(a=0\). For \(a\ge b\) bound the logarithm directly.
For \(0<a<b\), the left side increases with \(a\), since
\(\log(1+t)\ge t/(1+t)\) for \(t\ge0\). Apply this at \(a=U\).
Combine \eqref{trmean:budget-eq} with \eqref{trmean:lower} and
substitute \(U\le K_3+\mathbb E\ell_u\). The coefficient of
\(\mathbb E\ell_u\) is \(1+C_\phi\). The word counting bound and
\eqref{trmean:signed} now give
\[
 \mathbb E D_u\le K_0+(1+C_\phi)K_7+C_\phi(K_3+b)
           +\frac{(\eta_*+4\sqrt{e_*})H}{1-\delta_*}
           +(1+\kappa)L_P .
\]
The source budget adds the nonnegative term
\(\mathbb E\log|\{i:i\text{ is a word at }u\}|\) for nonempty word
sets. This is the only translate-dependent summand in its final error
bracket. We verify its numerical absorption below.
The two sufficient loss assumptions
in the source are \(\mathbb E(\mathcal A_u-G_u)\le L_P\) and
\(\mathbb E I_u\le L_P\). Multiplying the second by \(\kappa\ge0\)
and adding gives \eqref{trmean:signed}.

For comparison with the geometric hypotheses, the frame bound
\(\sum_i\|A_{ui}^*A_{ui}-I\|_{\rm HS}^2\le r_u/4\) gives
\(|M_u-r_u|\le r_u/2\) by Cauchy--Schwarz. The original target
minimum implies \(T_{\rm out}=0\), and integral columns with their
height budgets give \(\mathcal J=0\). For the source variation put
\[
\begin{aligned}
 V_0&=2mc_EB_0+c_H(n_Ta_E\mathcal L_a+mB_0)+Ac_{\rm count}\log k,\\
 K_3&=c_{\rm slope}n_Ta_E\mathcal L_a(2+m/k)+V_0+c_s mB_0 .
\end{aligned}
\]
The exponent, height and count comparisons give
\(|(E_{ui}+h_{ui})-(E_{ui'}+h_{ui'})|\le V_0\).
Indeed the exponent differences cost \(2mc_EB_0\), the differences
of \(h_{ui}+A\log T_{ui}\) cost the term with \(c_H\), and the
log-level differences cost \(Ac_{\rm count}\log k\).
Subtract the lower word-slope estimate in the preceding proof at
\(i'\) from the upper estimate
\(\lambda^{\max}_{ui}\le\lambda^{M,\max}_{ui}+G_u+F_u^{\,0}
-E_{ui}+c_smB_0\) at \(i\). Use
\(\lambda^{M,\max}_{ui}\le c_{\rm slope}n_Ta_E\mathcal L_a\)
and \(\mu^M_{ui'}\ge-C_M\).
This bounds \(\lambda^{\max}_{ui}-h_{ui}-
(\mu(S_{ui'})-h_{ui'})\) by \(K_3+\ell_{ui'}\).
Average first over \(i'\), maximize over \(i\), and take the
translate mean. This gives exactly the asserted bound on \(U\).

For the numerical absorption put \(y=x^\varepsilon\),
\(y_D=y^{e_{\rm div}},\ y_L=y^{e_{\rm loss}}\), and
\(P=n_Ta_E(1+m^2/k)+m(B_0+N_c)\). Here \(y,y_D,y_L\ge1\).
Let \(K\) bound the seven nonnegative coefficients
\(c_{\rm div},c_E,c_H,c_{\rm count},c_{\rm slope},c_s,c_{\rm loss}\).
The scalar size assumptions imply
\[
\begin{gathered}
 1\le n_T,\quad 1\le b\le a_E,\quad mb\le a_E,\quad
 m\le k,\quad 1\le k\le\min(a_E,x^{e_K}),\\
 A\le2n_Ta_E,\quad kb'\le a_E,\quad X''\le a_E,\quad
 0\le\mathcal L,\mathcal L_a\le2y,\quad B_0,N_c\ge0 .
\end{gathered}
\]
In particular \(mb\le a_E\) follows from \(m\le x^{(1+\varepsilon)/2}
\le x\) and \(xb\le a_E\). Also
\(\delta_*\le1/2,\ (\eta_*+4\sqrt{e_*})H\le Kyn_Ta_E\),
\(\log|\{i\}|\le(k-3)\log2\le a_E\) in each block, and
\(\log x\le y\). Put \(K'=1+K,\ h'=1+c_{\rm HH},\ g'=1+e_K\),
with \(c_{\rm HH},e_K\ge0\). From \(P\ge n_Ta_E\ge1\),
\(\log k\le e_K y\), and the displayed formulas one obtains
\[
\begin{aligned}
 K_0&\le12K y y_D P,&
 K_7&\le5(K')^2yP,&
 K_3+b&\le13K'g'yP,\\
 C_\phi&\le8h'y,& L_P&\le K y_L P.&
\end{aligned}
\]
For \(K_0\), the four terms cost at most
\(4KyP,(2y+y_D)KP,KP,4KyP\), respectively.
For \(K_7\), its term containing the word count is at most
\(4K^2y\,n_Ta_E m^2/k\), and \(c_sX''+1\le K'P\).
For \(K_3\) the base slope costs \(6KyP\), the variation costs
\((5+2e_K)KyP\), and \(c_s mB_0+b\le K'P\).
The estimate for \(C_\phi\) uses \(\log(1+x^{c_{\rm HH}})\le h'y\).
Thus, in units of \((K')^2h'g'y^2y_Dy_LP\), the five terms in the
displayed bound on \(\mathbb E D_u\) and the extra log word count
cost at most \(12,45,104,2,10,1\), respectively.
Their sum is \(174\). The common threshold \(y\ge200(K')^2h'g'\)
therefore gives \(\mathbb E D_u\le y^{3+e_{\rm div}+e_{\rm loss}}P\).
Substituting \(B_0\le2nmbx^\varepsilon\),
\(N_c\le n_2kb x^{2\varepsilon}/n^2\) and the specialized exponents
proves \eqref{scope:average}.
\end{proof}

\begin{corollary}[Room for lattice defects on the actual scales]\label{trmean:reserve}
Retain the arithmetic identifications, \eqref{trmean:lower}, the mass,
frame, variation and numerical bounds of Corollary~\ref{trmean:loss}.
Take \(e_{\rm div}\le2,\ e_{\rm loss}=0\). Allow arbitrary
\(\mathcal J,T_{\rm out}\) and write
\(\Delta=\mathbb E(\mathcal A_u-G_u+\kappa I_u)-(1+\kappa)L_P\).
At the same threshold \(\mathbb E D_u\le y^3P+
\Delta+4\mathcal J/(1-\delta_*)+T_{\rm out}\). In particular,
\begin{equation}\label{trmean:reserve-eq}
 \Delta+\frac{4\mathcal J}{1-\delta_*}+T_{\rm out}
                 \le (y^5-y^3)P
\end{equation}
suffices for \eqref{scope:average}. If \(\mathcal J=T_{\rm out}=0\)
and \(\Delta\le0\), the same window and character hypotheses yield
\(G_{\mathcal F_\psi}(x)\le x^{1/2+7\varepsilon}\) for large \(x\).
\end{corollary}
\begin{proof}
In the preceding proof replace \(12Kyy_DP\) by
\(12K\max(y,y_D)P\), using its four separate estimates. With
\(y_D\le y^2\) and \(y_L=1\), the same six coefficients now bound
the costs in units of \((K')^2h'g'y^2P\). The threshold absorbs
their sum into \(y^3P\). Keeping the actual signed loss in place
of its cap adds exactly \(\Delta\) to the upper bound.
Equation~\eqref{trmean:budget-eq} adds the two indicated nonnegative
terms. The variation bound still concerns the original source. Under
\eqref{trmean:reserve-eq} the resulting bound is \(y^5P\), and the
same bounds on \(B_0,N_c\) give \eqref{scope:average}. Under the last
assumptions use \(y^3P\) directly. Propositions~\ref{ar:mean} and
\ref{ar:four} with \(C=3\) give \(G_K\le x^{1/2+6\varepsilon}+\Sigma\).
The excluded twists satisfy \(\Sigma\le4o\,x^{1/2+\varepsilon}\).
Taking \(y\ge1+4o\) proves the stated window bound.
\end{proof}

The more general
Proposition~\ref{trmean:assembly} also allows the quantile bound
without a uniform lower bound on \(d_u/n_u\).
In both cases the order of selection remains that of
Section~\ref{scope:account}. For each prescribed late window
threshold one needs one complete window, then at every required
scale a single host with tied data and the stated mean bounds.
Proposition~\ref{ar:mean} then performs the four-twist extraction.
The reserve test \eqref{trmean:reserve-eq} gives the original window
exponent \(1/2+9\varepsilon\).
The group size incurs no extra logarithm. The uniform arithmetic
premises, including the recorded gains on the actual descended
columns, remain to be established on that family.
